% Heat-led computing in homes.
% Build: latexmk -pdf main.tex   (see README.md)
% Reproduction commands and result sources are collected in the appendix.
\documentclass[11pt,a4paper]{article}

\usepackage[T1]{fontenc}
\usepackage[utf8]{inputenc}
\IfFileExists{lmodern.sty}{\usepackage{lmodern}}{}
\usepackage[margin=2.5cm]{geometry}
\usepackage{amsmath,amssymb}
\usepackage{booktabs}
\usepackage{array}
\usepackage{siunitx}
\usepackage{graphicx}
\usepackage{pgfplots}
\pgfplotsset{compat=1.18}
\usepackage{xcolor}
\usepackage{hyperref}
\IfFileExists{xurl.sty}{\usepackage{xurl}}{} % break long URLs anywhere

\hypersetup{colorlinks=true, linkcolor=blue!50!black, citecolor=green!40!black,
  urlcolor=blue!60!black}

\sisetup{group-separator={,}, group-minimum-digits=4, group-digits=integer, detect-all}
\DeclareSIUnit\chf{CHF}

% Long file paths in \texttt: allow line breaks rather than overfull lines.
\sloppy
\setlength{\emergencystretch}{3em}

% Keep tables visually separated from surrounding prose and from their captions.
% The slightly larger values are intentional: the paper contains long captions and dense tables.
\setlength{\textfloatsep}{24pt plus 4pt minus 2pt}
\setlength{\intextsep}{20pt plus 4pt minus 2pt}
\setlength{\floatsep}{16pt plus 3pt minus 2pt}
\setlength{\belowcaptionskip}{8pt}

% Visible markers for unfinished content.
% \todo{...}    : text that still has to be written (skeleton section, missing source).
% \pending{...} : a result that is being computed right now (LP experiment runs).
\newcommand{\todo}[1]{\textcolor{red!75!black}{\textbf{[TODO:~#1]}}}
\newcommand{\pending}[1]{\textcolor{orange!85!black}{\textbf{[pending]}\if\relax\detokenize{#1}\relax\else~#1\fi}}

% Notation
\newcommand{\Sites}{\mathcal{S}}
\newcommand{\Win}{\mathcal{T}_\tau}
\newcommand{\Coh}{\mathcal{C}}

\title{Heat-led computing in homes:\\
  scheduling, sizing and economics of servers that heat a house}
% Author order rotates among the team; this paper starts with Ennio.
\author{Ennio Binder \and Jona Richartz \and Nikita Sosnovski}
\date{22 November 2026}

\begin{document}
\maketitle

\begin{abstract}
Heat-led computing promises two services from one electrical input, but heat demand and paid
computing demand need not coincide. We simulate the case's Basel reference house and labelled
variants with an event-driven fleet model that draws finite job demand from sourced volumes and
traces, schedules it with one shared heat-first policy, and settles every house with the case's
lossless heat balance. With unlimited work, four abstract \qty{1}{kW} servers could cover 58.0\% of
the house's heat at 64.7\% utilisation; with the sourced demand (scaled to one house), they cover
22.8\% at 25.4\% utilisation, dump no heat and miss 8.9\% of job deadlines, mostly in summer when
only hot water is drawn. Income exceeds the electricity cost, but four GPU servers lose about
CHF~14,300 a year after hardware; one server earns almost the same income, and the combined system
of operator and host keeps about CHF~5,250 a year (the operator alone CHF~4,470, payback 3.4 years). Calibrated to an organiser-provided real house and assuming 86\% boiler efficiency,
the case formula's fitted design heating power (\qty{10.13}{kW}) is within 1.3\% of the case's \qty{10}{kW}, while summer heat use is three
times the case's hot water and absorbs nearly all summer work. Routing to a year-round heat sink
moves 76\% of July work to an indoor pool in a local scenario, and in synthetic 33-site experiments
with earlier research strategies it reduces missed deadlines where houses alone cannot. A demonstrator shows the same scheduler
freezing and resuming a real transcription and language-model container on the heat signal. These
results state measurable conditions for a viable installation; they do not claim realised
profitability.
\end{abstract}

\section{Introduction}
\label{sec:intro}

A server converts most of its electrical input into heat. In a home, that heat can displace
heat-pump or backup-heater operation, while the server sells computing work. The appeal is simple:
one electricity bill may provide both a computing service and useful heat. The hard question is
whether the two uses coincide often enough, and whether the value created covers the equipment and
operation costs for both the operator and the host.

This question is especially sharp for a single-family house. In the Basel reference case, hourly
weather in 2025 implies a peak demand equivalent to 12.5 servers of 1~kW, while summer hot water
uses only about 0.6 servers. A system sized for the winter peak is therefore poorly utilised in
summer; a small system leaves much of the winter demand to backup heat. Scheduling can shift work
within those physical limits, but cannot create compute customers or heat demand. The business
case must consequently be tested against both the heat balance and the price of sellable work.

We study the following questions:
\begin{description}
  \item[Q1] What heat coverage, utilisation and deadline performance can a case-compliant single
    house achieve, first with unlimited work and then with finite, sourced job demand?
  \item[Q2] Under which conditions does the single-house system cover its hardware, and how many
    servers per house does that favour?
  \item[Q3] In a multi-site setting, can moving deadline-bound work to sites with year-round heat
    demand reduce missed deadlines?
\end{description}

\paragraph{Approach.} We use the case's fixed Basel reference house and hourly weather, and add
labelled variants: a house calibrated to an organiser-provided gas record, a local scenario with an
indoor pool and a synthetic Swiss network. An event-driven fleet model draws jobs from sourced
volumes and traces, quotes prices and schedules work with one heat-first policy shared by the
simulation and the real-container controller. A rolling-horizon linear programme and a full-year
upper bound serve as analytical comparisons. For economics, we separate the case break-even
(electricity net of heat value) from full cost after hardware, and expose the volume, price and
server count at which the house breaks even. Network results are scenario analyses, not forecasts
of a deployed market.

\paragraph{Contributions.} This paper makes four claims, each limited to the evidence available.
\begin{enumerate}
  \item A reproducible reference-house analysis that separates the physical limit (unlimited work)
    from the result under sourced finite demand, and locates the full-cost boundary in volume,
    price and server count (Sections~\ref{sec:case}, \ref{sec:economics} and~\ref{sec:evaluation}).
  \item One heat-first scheduling policy used by the fleet simulation and a real-container controller,
    with greedy, network-greedy and linear-programme strategies as research comparisons
    (Sections~\ref{sec:model} and~\ref{sec:scheduling}).
  \item A check of the case's house model against a real house's measured gas use, and network
    scenarios showing when routing to year-round heat sinks changes deadline outcomes
    (Section~\ref{sec:evaluation}).
  \item A demonstrator in which the same scheduler freezes and resumes a real transcription and
    language-model container on the heat signal, with measured CPU time and energy
    (Section~\ref{sec:demonstrator}).
\end{enumerate}

The paper's practical target is not to claim that any server makes a house cheaper to heat. It is
to identify a credible single-house path: enough paid work at a sufficient realised price, a
hardware and installation cost compatible with that revenue, and heat that displaces an actual
alternative. If those conditions cannot be met, the result should make that failure boundary
legible to a homeowner and to DigitalHeat.

\section{Related work}
\label{sec:related}

\paragraph{Existing installations.} Heat reuse from computing is already demonstrated at several
scales. Centrica describes a trial in which Heata attaches domestic servers to hot-water cylinders,
pays their electricity and processes the partner's workloads \cite{heata_centrica2025}. In Basel,
IWB reported a Neutral-IT server-heating pilot in a residential building with conventional
heating retained as backup \cite{iwb_neutralit}. Infomaniak reports a Geneva data centre feeding
district heating \cite{infomaniak_d4}. These primary accounts establish examples of integration,
but do not provide a comparable public dataset of utilisation, operator receipts, capital cost or
failure rates. We therefore use them as context, not evidence of profitability for our proposed
home node.

\paragraph{Academic context.} Antal and colleagues already studied workload scheduling for
servers placed in homes as a primary heat source. Their constraint-satisfaction approach
targets room-temperature set points and workload deadlines, using a thermodynamic and learned
heat model validated against monitored distributed equipment \cite{antal2021}. Thus neither
residential placement nor thermal-aware scheduling alone is novel here. Our narrower question
uses the case's water-based heat balance and annual hourly demand, compares a greedy rule with
a rolling-horizon divisible-work LP, and separates operator cash flow from the host's displaced
heating expense. District-heating studies such as Wahlroos and colleagues' two-actor analysis
address a larger, more stable heat sink \cite{wahlroos2017}; our house-scale seasonality and
customer-fill constraints are different. We do not infer that either earlier result validates
our unmeasured hardware or commercial assumptions.

\paragraph{Data used.} Hourly weather comes from Open-Meteo \cite{openmeteo}, the electricity price
from ElCom \cite{elcom}. The case points to MLPerf Inference for energy per job \cite{mlperf} and to
the Alibaba GPU cluster trace as inspiration for job durations \cite{alibaba_trace}.

\section{Case and reference house}
\label{sec:case}

All reference-case figures use the house exactly as defined in the challenge specification;
other houses and sites appear only as labelled variants (Appendix~\ref{app:repro}).

\paragraph{Weather and heat demand.} Hourly air temperature $T$ for Basel (47.56\,N, 7.59\,E),
1~January to 31~December 2025, local time Europe/Zurich, comes from the Open-Meteo historical
weather API \cite{openmeteo}. A heating day is a local calendar day whose mean hourly temperature is
below \qty{12}{\celsius}. On heating days the space-heating demand per hour is
\begin{equation}
  D^{\mathrm{space}}_t = \max\!\left(0,\ \qty{10}{kW} \cdot \frac{20 - T_t}{25.6}\right)\cdot \qty{1}{h},
  \label{eq:heating}
\end{equation}
where $25.6 = 20 - (-5.6)$ and \qty{-5.6}{\celsius} is the coldest hour of 2025. Hot water adds
\qty{12}{kWh} per day, all year. With the 2025 weather this gives 206 heating days,
\qty{26850.55}{kWh} of space heat and \qty{4380}{kWh} of hot water, i.e.\ \qty{31230.55}{kWh} of heat
demand. The calculation and input files are listed in Appendix~\ref{app:repro}.

\paragraph{Storage, server and prices.} The house has one lossless tank with \qty{35}{kWh} usable
capacity. A server draws \qty{1}{kW} at full load, of which \qty{80}{\percent} is recoverable, i.e.\
\qty{0.8}{kWh} of heat per server-hour (\qty{19.2}{kWh} per server-day). Electricity costs
CHF~0.3325 per kWh (ElCom 2026, IWB Basel, category H4, one flat price \cite{elcom}). Used server
heat is valued at the price of the heat pump it replaces, electricity price $\div$ 2.5, i.e.\
CHF~0.133 per kWh of heat. Heat that the servers do not cover comes from a backup heater.

\paragraph{The central tension.} The coldest hour needs $\qty{10}{kW} \div \qty{0.8}{kW} \approx 12.5$
servers; summer hot water needs $\qty{12}{kWh} \div \qty{19.2}{kWh} \approx 0.6$ servers
(\texttt{docs/case/key-figures.md}).

\section{System model}
\label{sec:model}

This section states the model compactly; the full parameter tables with units and sources are in
\texttt{docs/model/routing-optimisation.md} (decision 0014), and the implemented variant in
\texttt{docs/model/network-scheduler.md} (decisions 0024, 0025).

\paragraph{Sites and hours.} A network consists of sites $s \in \Sites$ (houses, hotels with spa,
indoor pools). Time runs in hourly steps $t$, local time Europe/Zurich. Site $s$ has $N_s$ servers,
recovers $a_s$ kWh of heat per server-hour ($a_s = 0.8$ in the case), has a minimum load share
$\mu_s$ while on, and a lossless tank of usable capacity $B_s$ kWh. Its gross heat demand (space
heating and hot water) in hour $t$ is $D_{s,t}$ kWh.

\paragraph{Realised heat balance.} Whatever a strategy decides, every hour is settled by one shared
function (\texttt{HeatBalance.Step}) on the \emph{actual} demand: server heat $a_s u_{s,t}$ first
covers demand, then charges the tank, and only the rest is dumped; the tank discharges for the
remaining demand, and the backup heater covers what is still missing. There is no second heat
model: the year simulation, the network simulation and the tests share this function.

\paragraph{Work.} We use two job types (decision 0020). \emph{Type~1} is paid batch LLM work
(for example summarising public documents); it arrives in \emph{cohorts} $c \in \Coh$, one per
origin site $o_c$ and day, at local 00:00, with an exclusive deadline $d_c$ \qty{24}{h} later and a
remaining amount $W_c$ server-hours. Its volume in our experiments is a labelled scenario: a share
of the origin's yearly heat-led server-hours, because no sourced arrival rate exists. Type~1 may
run only at eligible sites $E_c \subseteq \Sites$, which the data-locality rules (decision 0013)
determine before any optimisation; no price or heat term can override them. \emph{Type~2} is
\emph{filler}: research simulation with checkpoints and no deadline, available up to
$F_{s,t} = \varphi N_s$ server-hours per hour. $\varphi = 1$ is the case-permitted ``unlimited
jobs''; lower values model a shortage of sellable work. Work is divisible into small units, and
servers can run at part load, so loads are continuous.

\paragraph{Forecasts.} Plans for future hours use a demand forecast $\hat D_{s,t}$. By default it is
perfect foresight (the actual 2025 demand) and labelled as such. A synthetic forecast error follows
an AR(1) process in the lead time $k$:
\begin{equation}
  e_k = \rho\, e_{k-1} + \sigma_k \xi_k,\qquad \sigma_k = \sigma_{24}\sqrt{k/24},\qquad
  \xi_k \sim \mathcal{N}(0,1),\quad e_0 = 0,
  \label{eq:ar1}
\end{equation}
applied as a relative error to the demand and seeded per site and hour, so runs are reproducible.
The values of $\sigma_{24}$ and $\rho$ are \texttt{TODO(source)}; we therefore show a sweep
($\rho = 0.8$ is a placeholder).

\paragraph{Planning model.} At hour $\tau$ the scheduler plans the window
$\Win = \{\tau, \dots, \tau + H\}$ with $H = 48$ and commits only hour $\tau$ (rolling horizon,
decision 0014 \S6.2). It knows the realised tank content $S^0_s$ and the cohorts that have arrived
but are not finished; it does not know future arrivals. All variables are continuous and
non-negative: total load $u_{s,t}$, Type~1 work $x_{c,s,t}$ (defined only for $s \in E_c$ and
$t < d_c$), filler $f_{s,t}$, heat used directly $q^{\mathrm{dir}}_{s,t}$, tank charge
$q^{\mathrm{ch}}_{s,t}$ and discharge $q^{\mathrm{dis}}_{s,t}$, backup heat $b_{s,t}$, dumped heat
$z_{s,t}$, tank content $S_{s,t}$ and missed work $m_c$. For all $s \in \Sites$, $t \in \Win$:
\begin{align}
  u_{s,t} &= \textstyle\sum_{c} x_{c,s,t} + f_{s,t}, \qquad
  \mu_s N_s \le u_{s,t} \le N_s, \qquad f_{s,t} \le \varphi N_s, \label{eq:load}\\
  a_s u_{s,t} &= q^{\mathrm{dir}}_{s,t} + q^{\mathrm{ch}}_{s,t} + z_{s,t}, \label{eq:produce}\\
  q^{\mathrm{dir}}_{s,t} + q^{\mathrm{dis}}_{s,t} + b_{s,t} &= \hat D_{s,t}, \label{eq:demand}\\
  S_{s,t+1} &= S_{s,t} + q^{\mathrm{ch}}_{s,t} - q^{\mathrm{dis}}_{s,t}, \qquad
  0 \le S_{s,t} \le B_s, \qquad S_{s,\tau} = S^0_s, \label{eq:tank}\\
  \textstyle\sum_{s,t} x_{c,s,t} + m_c &= W_c \quad (d_c \le \tau + H), \qquad
  \textstyle\sum_{s,t} x_{c,s,t} \le W_c \quad (d_c > \tau + H). \label{eq:work}
\end{align}
For the current hour $\hat D_{s,\tau}$ is the measured demand. Eligibility, deadlines and
locality are enforced by which $x_{c,s,t}$ exist, not by penalties. Work processed at $s \ne o_c$
counts as \emph{moved}. A uniform transfer delay $\delta$ (whole hours, default 0; decision 0030)
models friction under the assumption ``data staged at arrival'': the cohort's data starts copying to
the eligible sites when it arrives, so $x_{c,s,t}$ with $s \ne o_c$ exists only for
$\alpha_c + \delta \le t < d_c$, where $\alpha_c$ is the arrival hour. This is a lower bound on
transfer friction; $\delta \ge \qty{24}{h}$ means no moves. There is no distance model and no
transfer cost (\texttt{TODO(source)}: a sourced data volume per server-hour); runs report moved
server-hours instead. The objective is lexicographic and is given in Section~\ref{sec:scheduling}.

\paragraph{Why an LP.} The problem is a flow in a time-expanded network with nodes $(s,t)$: work
waits at a site until its deadline, moves to an eligible site, or is computed; heat is a second
commodity that is used, stored or dumped, coupled to work by \eqref{eq:produce}
(\texttt{docs/model/routing-optimisation.md}~\S6.1). With divisible work and part load it stays a
linear programme. Integer servers, indivisible or non-preemptive jobs, start-up costs or
all-or-nothing payments make it a MILP (future work, Section~\ref{sec:limits}).

\section{Scheduling}
\label{sec:scheduling}

Every strategy answers the same question each hour: how much should each site's servers run, and
which work should they do? All strategies share one heat balance per site (the case's lossless
tank, \qty{80}{\percent} recovery, backup heater for the rest) and one actuation path, so the
scheduler that computes the year is the one that drives a real container in the demonstrator
(Section~\ref{sec:demonstrator}).

\subsection{Five strategies}

We compare strategies of increasing intelligence.

\textbf{B0 (always on)} runs every server all the time, $u_{s,t} = N_s$, and dumps what the house
cannot take.

\textbf{B1 (local greedy, heat-led)} runs exactly what the house needs now plus what fits into the
tank, and never dumps:
\begin{equation}
  u_{s,t} = \operatorname{clamp}\!\left(\frac{D_{s,t} + B_s - S_{s,t}}{a_s},\ \mu_s N_s,\ N_s\right).
  \label{eq:greedy}
\end{equation}
This is the scheduler that produces the scored figures of the reference house. With jobs, each site
first serves its own Type~1 work in earliest-deadline-first order and then fills the rest with
filler; it never moves work.

\textbf{B2 (network greedy)} keeps B1's load at every site but moves paid work (Type~1, with a
24-hour deadline) that its home site cannot finish in time to the eligible site with the most free
heat-led capacity. Whether the origin can finish in time is estimated by persistence. With a
transfer delay $\delta$, a cohort may move only from hour $\alpha_c + \delta$ on.

\textbf{LP (rolling horizon)} plans all sites over a rolling 48-hour window
(Section~\ref{sec:model}) and decides, in a strict order of priorities that mirrors the case's
ranking, how to minimise dumped heat, maximise used heat, avoid missed deadlines, maximise
server-hours and, last, keep work at home unless moving helps.

\textbf{UB (full-year LP)} solves the same model once for the whole year with perfect knowledge of
all demand. It is not a strategy that could run; it serves only as an upper bound.

The LP is a linear programme because our work splits into small units (documents, simulation
tasks) and servers can run at part load. This keeps it fast: a full year of rolling-horizon solves
on the reference house takes about \qty{45}{s}, and on 33 sites about \qty{0.22}{s} per hour in
January and \qty{0.38}{s} per hour in July (Debug build, one core), i.e.\ roughly 30 to 60 minutes
per year (\texttt{docs/model/network-scheduler.md}). Integer servers or non-splittable jobs would
make it a MILP (future work). Weather-based demand forecasts are perfect foresight unless a forecast
error \eqref{eq:ar1} is added; we label every result accordingly.

\subsection{Lexicographic objective}

The LP solves a sequence of stages over the window. Let $J_1, \dots, J_5$ be
\begin{align*}
  J_1 &= \textstyle\sum_{s,t} z_{s,t} \ (\min), &
  J_2 &= \textstyle\sum_{s,t} \bigl(q^{\mathrm{dir}}_{s,t} + q^{\mathrm{dis}}_{s,t}\bigr) \ (\max), &
  J_3 &= \textstyle\sum_{c} m_c \ (\min),\\
  J_4 &= \textstyle\sum_{s,t} u_{s,t} \ (\max), &
  J_5 &= \textstyle\sum_{c} \sum_{s \ne o_c} \sum_{t} x_{c,s,t} \ (\min). &&
\end{align*}
Stage $k$ optimises $J_k$ subject to \eqref{eq:load}--\eqref{eq:work} and to every earlier stage
staying at its optimum $J_i^\star$ within a tolerance of $10^{-9}$ (relative) plus $10^{-9}$
(absolute), $i < k$. If $J_1^\star = 0$, dump is fixed to exactly zero (heat-led ``hard'' mode).
On one site without jobs only stages $J_1$, $J_2$ and $J_4$ apply. The order encodes a choice: heat
comes first, money only through the server-hours that heat allows. Section~\ref{sec:economics}
shows what this choice costs.

\subsection{The greedy rule as a special case}

Take one site, a one-hour window, no Type~1 work and unlimited filler. In that LP an extra
server-hour is worth $p^\varphi - eP + v a$ while it covers demand, $p^\varphi - eP + \theta v a$
while it charges the tank (with $\theta$ the value of heat left in the tank at the window end) and
$p^\varphi - eP - \pi^{\mathrm{dump}} a$ when it is dumped, where $p^\varphi$ is the filler price,
$e$ the electricity price, $P$ the server power, $v$ the value of used heat and $\pi^{\mathrm{dump}}$
a dump penalty. If charging pays and dumping does not, the LP runs exactly until the tank is full,
which is \eqref{eq:greedy} (\texttt{docs/model/routing-optimisation.md}~\S8.1). The LP therefore
generalises the greedy rule in three directions: \emph{time} (a forecast, $H > 1$), \emph{space}
(moves, $|\Sites| > 1$) and \emph{markets} (paid work versus filler).

\subsection{Local timing versus the network}

Part of what moving achieves could also be achieved locally: in summer a site's heat-led budget
(hot-water draw plus free tank room) is scarce, and B1 spends free tank room on filler as soon as
the day's Type~1 work is done. Two control strategies (decision 0029) separate \emph{local timing}
from \emph{the network}. \textbf{B1-reserve} and \textbf{B2-reserve} work like B1 and B2, but filler
may not fill the tank beyond a reserve kept free for tomorrow's paid work:
\begin{equation}
  R_s = \min\!\Bigl(B_s,\ a_s \cdot \max\bigl(0,\ \hat W_s - \bar D_s\, d / a_s\bigr)\Bigr),
  \label{eq:reserve}
\end{equation}
where $\hat W_s$ is the Type~1 work that arrived at the site in the last \qty{24}{h} (persistence
estimate of the next cohort), $\bar D_s$ the mean demand of the last 24 observed hours and
$d = \qty{24}{h}$ the deadline. Filler in the current hour is capped so that the tank keeps at least
$R_s$ free. Sites without Type~1 behave exactly like B1. The reserves use no forecast.

\section{Economics: who receives the heat value?}
\label{sec:economics}

The case break-even is an accounting measure for the combined operator and host. For one
full-load-equivalent hour, the case server consumes \qty{1}{kWh} of electricity at
CHF~0.3325/kWh and delivers \qty{0.8}{kWh} of useful heat worth
CHF~$0.3325/2.5=0.133$ per kWh. If all that heat displaces heat-pump operation, the combined
electricity cost net of heat value is $0.3325-0.8(0.133)=\text{CHF}~0.2261$ per hour.
For the four-server fleet run of the reference house, the same calculation is
$(\text{CHF}~2{,}964.84-\text{CHF}~948.75)/8{,}916.80
=\text{CHF}~0.2261$ per sold server-hour
(Appendix~\ref{app:repro}). This number excludes equipment and operation.

\paragraph{Operator and host.} In the provisional operator model, DigitalHeat
owns the hardware, sells the computing and pays or reimburses its electricity. The homeowner
receives the useful heat and retains a backup heat source. With free heat, the host's potential
avoided heat-pump expense is $Q_{\mathrm{use}}v$, while DigitalHeat's cash balance is computing
receipts minus electricity, depreciation, operations, wear and installation. The CHF~949 heat
value in the four-server run belongs to the host; it is not also operator income. If the host pays
a fee $h$ per useful kWh, $Q_{\mathrm{use}}h$ transfers from host to operator without changing
combined surplus. A positive combined surplus alone does not establish a viable offer for both.

\paragraph{Full cost under sourced demand.} With the sourced job demand of the fleet model, the
four-server reference house earns CHF~13,360 a year (CHF~8,911 from Type~1, CHF~4,449 from
Type~2). Electricity costs CHF~2,965, so the margin, which includes the heat value, is CHF~11,344. The
L4 hardware profile, normalised to a \qty{1}{kW} server at CHF~34,244 over a six-year accounting
life plus CHF~708 yearly operations, costs CHF~6,415 per server-year, or CHF~25,661 for four
servers. The net after hardware is $-\text{CHF}~14{,}317$ per year. Table~\ref{tab:breakeven}
changes one input at a time to find where the net reaches zero.

\begin{table}[ht]
  \centering
  \caption{Hardware break-even of the reference house in the fleet model (Basel 2025, L4 profile,
    sourced demand scaled by 0.1 unless stated). Net = income + heat value $-$ electricity $-$
    hardware per year.}
  \label{tab:breakeven}
  \begin{tabular}{>{\raggedright\arraybackslash}p{0.38\linewidth}>{\raggedright\arraybackslash}p{0.54\linewidth}}
    \toprule
    Change & Result \\
    \midrule
    Servers per house & 1: $+$CHF~5,252, payback 3.1 years, 18.9\% coverage, 5.9\% missed deadlines;
      2: $-$CHF~621; 3: $-$CHF~7,654; 4: $-$CHF~14,317 \\
    Job volume (all types), 4 servers & break-even at 2.8 times the scaled demand; 30\% of jobs then
      miss their deadline, and above about 3 times the margin levels off near CHF~35,000 because the
      house cannot use more heat \\
    Type~1 prices, 4 servers & break-even at 2.61 times the sourced prices (LLM inference about
      CHF~5.81 per L4 server-hour instead of CHF~2.23) \\
    Type~2 price band only & no break-even: with elasticity $-0.8$, the accepted volume falls almost
      as fast as the price rises, so revenue above the competitor price grows only with $p^{0.2}$ \\
    \bottomrule
  \end{tabular}
\end{table}

At the scaled demand, income hardly depends on the number of servers (CHF~13,300 to 14,100 for one
to four), because work is limited by demand in winter and by heat in summer, not by computing
capacity. The hardware therefore decides the result, and the single-server house is the only
configuration that covers it at this demand. All nets credit the heat value to the combined system; for the operator alone,
who gives the heat away, they are lower by that value (CHF~949 with four servers, CHF~784 with one,
so the single-server operator keeps about CHF~4,470). The price route is optimistic: Type~1
customers do not react to price in the model. The volume baseline is itself a scenario value (Section~\ref{sec:limits}).

\paragraph{Upper bound with unlimited work.} Two L4 purchase costs appear in this section: the
hardware profile's CHF~34,244 per normalised \qty{1}{kW} server, used by the fleet results above,
and a later card-price sensitivity of CHF~44,179. Neither is a quoted installable chassis. If every
heat-feasible hour were sold (the case's unlimited-work assumption, \num{22661} server-hours for four
servers), the four-server, free-heat operator would need an average realised compute price of
\begin{equation}
  p_{\mathrm{op}}^{*}
  =0.3325+\frac{4(44{,}179/6+708)}{22{,}661}
  \simeq \text{CHF}~1.76\text{/server-hour}
  \label{eq:operator-break-even}
\end{equation}
at the later cost, or about CHF~1.46 at the profile's cost, with six-year depreciation and CHF~708
yearly operations per server. The calculation excludes installation, cooling, financing, platform
fees, idle power and wear. Consumer-card deployment is excluded from the proposed architecture
pending licence clarification; the cost ledger is identified in Appendix~\ref{app:repro}.

Two token-price bases give different hourly revenue. The fleet job type uses \num{25} million tokens
per L4 server-hour and the Public AI Apertus~8B price, about CHF~2.23 per server-hour. An earlier
benchmark for another 8B model (\num{14.36} million tokens per server-hour) gives CHF~1.81 at that
competitor's list price and about CHF~0.90 at its batch price. Neither is a DigitalHeat sale: both
transfer benchmarks for a different model and GPU, and no recurring Type~1 order is documented. In
the unlimited-work scenario at the later purchase price, with all hours sold at that list price,
four servers would leave about CHF~1,172 per year before wear and installation; this is not a fleet
result. Wear is not included: the available wear estimate was derived for a different card count
and utilisation and cannot be applied per sold hour here. The CPU-only measured research kernel is
also not evidence of a full-power, billable GPU filler workload.

\paragraph{Market comparison.} The dated price comparison in Appendix~\ref{app:repro}
provides a screening test, not a sales
forecast. At the profile's L4 hardware cost, full heat credit and the unlimited-work \qty{64.7}{\percent} assumed
sold hours, it estimates CHF~0.054 per million 8B-class tokens; at \qty{35}{\percent} sold
hours (above the fleet model's 25.4\%, Section~\ref{sec:eval-house}), CHF~0.093. A listed Swiss 8B offer was CHF~0.034 per million tokens on the same
date. The estimates use benchmark throughput for a different GPU workload, compare models
with unequal quality and omit the operator/host split. Under free heat and the later L4
purchase sensitivity above, they cannot establish a profitable or price-leading offer.

\paragraph{A smaller residential candidate.} A separate exploratory variant sizes a physical
node below the home's year-round hot-water demand. Four L4 cards on an R570-class host respect
Dell's documented four-card limit; the existing host estimate plus card TDP gives about
\qty{0.4505}{kW} electrical power and, at the case's 80\% recovery assumption,
\qty{0.3604}{kW} recoverable heat. This is below the case's constant
$12/24=\qty{0.5}{kW}$ hot-water demand. At continuous operation the variant could deliver at
most \qty{3157}{kWh} useful heat, worth about CHF~420/year to the host at the case's heat-pump
comparison. These are estimated physical quantities, not measured wall power or installed heat
recovery (Appendix~\ref{app:repro}). The fixed reference house and
its abstract \qty{1}{kW} server remain unchanged.

At a retail comparison price of CHF~0.2924 per GPU-hour, four billable GPUs, the estimated
power and a six-year accounting life, that exploratory node needs approximately 5,376 paid
full-node hours per year under the 9\% repair stress scenario, before missing integration and
service costs. This is about 61\% of the year. The retail price is not an achievable DigitalHeat
net contract price; paid reservations can exist without active heat production. The proposed
experiment is therefore a measured, thermally integrated small node coupled to a minimum-volume
customer offer, not a demonstrated profitable installation. Its operator and host cash flows must
be recalculated from measured whole-node power, useful heat and net contracted receipts.

\section{Demonstrator}
\label{sec:demonstrator}

\textit{Skeleton.}

The same scheduler code that produces the yearly balance also controls real running jobs (case
must-have). Each hour, \texttt{IScheduler.Decide(HeatSignal)} returns a target load per server;
\texttt{SchedulerController.Apply} passes it to an \texttt{IExecutor}, which is either simulated or
a \texttt{DockerJobExecutor} that pauses and resumes a container (decisions 0005 and 0009). In the
network simulation, an adapter feeds each site's committed decision through this unchanged path and
rejects any signal that differs from the one the decision was made on.

\todo{measured workloads: whisper.cpp transcription (\texttt{docs/jobs/transcription.md}),
Apertus 8B summaries (\texttt{docs/jobs/llm-summary.md}), research simulation (T-30)}

\todo{measurement method and scaling to the abstract 1~kW server; limits (CPU package versus wall
power)}

\todo{figure F4: January and July day with demand, tank, server load and the real job paused
(\texttt{DigitalHeat.LiveRun} replay output)}

\section{Evaluation}
\label{sec:evaluation}

We report the case's metrics: coverage (share of the total heat demand covered by server heat; the
rest comes from the backup heater), dumped heat, utilisation (server running hours $\div$ (servers
$\times$ 8{,}760)), missed deadlines in per cent of all jobs (fleet runs) or of Type~1 jobs (network experiments), and the work moved between sites.
Commands and source files are listed in Appendix~\ref{app:repro}; network CSV rows also
record their exact command lines.

\subsection{Reference house}
\label{sec:eval-house}

Table~\ref{tab:house-fleet} separates what the house could absorb from what it absorbs under
sourced demand. With unlimited work, four servers cover \qty{58.05}{\percent} of the heat demand at
\qty{64.7}{\percent} utilisation without dumping. This is the physical limit of the house; the
fleet model reproduces it within \qty{0.02}{\percent} (\num{22657.0} against \num{22661.0}
server-hours) when it is given a backlog that never runs out. With the sourced job demand, scaled
to one house, the same four servers run \num{8917} server-hours: coverage falls to
\qty{22.8}{\percent} and utilisation to \qty{25.4}{\percent}, still without dumping. Of
\num{77840} jobs, 6,905 (\qty{8.9}{\percent}) miss their deadline, and \num{5763} of the
\num{5771} jobs that never finished arrived between May and September: with only hot water to heat, the
scheduler has no useful heat to run them on. Four servers is the case's comparison count, not an
optimised one (Section~\ref{sec:eval-sizing}).

\begin{table}[ht]
  \centering
  \caption{Reference house and real house in the fleet model (Basel 2025, L4 hardware profile). The
    fleet columns use the same jobs and seed; the unlimited-work column is the case's physical
    limit from the year model (cross-checked by the fleet model).}
  \label{tab:house-fleet}
  \small
  \setlength{\tabcolsep}{4pt}
  \begin{tabular}{lrrrr}
    \toprule
    & \multicolumn{3}{c}{Reference house} & Real house \\
    \cmidrule(lr){2-4}\cmidrule(lr){5-5}
    & unlimited work & 4 servers & 1 server & 4 servers \\
    \midrule
    Heat demand [kWh] & \num{31231} & \num{31231} & \num{31231} & \num{41356} \\
    Heat demand covered & \qty{58.05}{\percent} & \qty{22.8}{\percent} & \qty{18.9}{\percent} & \qty{21.6}{\percent} \\
    Heat dumped [kWh] & 0 & 0 & 0 & 0 \\
    Utilisation & \qty{64.7}{\percent} & \qty{25.4}{\percent} & \qty{84.1}{\percent} & \qty{31.9}{\percent} \\
    Server-hours sold & \num{22661} & \num{8917} & \num{7366} & \num{11164} \\
    Missed deadlines & -- & \qty{8.9}{\percent} & \qty{5.9}{\percent} & \qty{0.01}{\percent} \\
    Computing income [CHF] & -- & \num{13360} & \num{13333} & \num{13780} \\
    Net after hardware [CHF/year] & -- & \num{-14317} & \num{5252} & \num{-14406} \\
    Case break-even [CHF/server-h] & 0.2261 & 0.2261 & 0.2261 & 0.2261 \\
    \bottomrule
  \end{tabular}
\end{table}

\paragraph{Real house.} The organisers provided the monthly gas readings of a real, anonymised
house that heats with gas. We convert January to September 2025 (later months are missing) to heat
with the 86\% minimum seasonal efficiency of new gas boilers \cite{eu813}, and fit the case
formula's two house parameters by least squares on those nine months. The fit gives
\qty{10.13}{kW} design heating power, against the case's \qty{10}{kW}, and \qty{38.8}{kWh} of heat
per day outside the heating season, against the case's \qty{12}{kWh} of hot water ($R^2 = 0.996$,
monthly errors $-5.9$ to $+8.4$\%). The case formula therefore describes this house's space heating
well; its summer use, which includes circulation and boiler losses that monthly data cannot
separate, is three times the case's hot water. Run with the same jobs, the real house misses 11
deadlines instead of 6,905 (Table~\ref{tab:house-fleet}): its summer heat absorbs the summer work.
Its income rises only slightly, because the extra work is mostly low-priced Type~2. The boiler's
real efficiency and age are unknown; the 86\% is an upper value for a new appliance.

\paragraph{Planning on one house.} The following comparison uses the network model with unlimited
work (commit \texttt{1863f57}). On a single house the greedy rule B1 is optimal: it reaches the full-year upper bound exactly, and
the rolling-horizon LP cannot improve on it (Table~\ref{tab:house-strategies}). The reason is that on
the reference house the minimum load is zero, the tank is lossless and work is unlimited, so using
heat now never reduces what can be used later; four servers are the limit in cold hours whatever the
strategy. Always-on (B0) reaches the same coverage but only adds \qty{9903}{kWh} of dumped heat in
mild hours. Planning on imperfect forecasts costs little: coverage falls from \qty{58.05}{\percent}
to \qty{57.82}{\percent} at $\sigma_{24} = \qty{20}{\percent}$, because the LP plans the current hour
for the forecast world and its plan can be slightly worse in reality.
\emph{On one house, intelligence does not pay; heat physics decides.}

\begin{table}[ht]
  \centering
  \caption{Reference-house strategies (four servers, \qty{35}{kWh} tank).
    Forecast errors are synthetic AR(1) innovations ($\rho=0.8$); UB is a bound.}
  \label{tab:house-strategies}
  \small
  \setlength{\tabcolsep}{4pt}
  \begin{tabular}{llrrrr}
    \toprule
    Strategy & Information & Coverage & Dumped [kWh] & Server-hours & Utilisation \\
    \midrule
    B0 always on & current hour & \qty{58.05}{\percent} & \num{9903} & \num{35040} & \qty{100}{\percent} \\
    B1 local greedy & current hour & \qty{58.05}{\percent} & 0 & \num{22661.0} & \qty{64.67}{\percent} \\
    LP rolling horizon & perfect foresight & \qty{58.05}{\percent} & 0 & \num{22661.0} & \qty{64.67}{\percent} \\
    LP rolling horizon & error $\sigma_{24} = \qty{5}{\percent}$ & \qty{58.03}{\percent} & 0 & \num{22654.7} & \qty{64.65}{\percent} \\
    LP rolling horizon & error $\sigma_{24} = \qty{10}{\percent}$ & \qty{57.99}{\percent} & 0 & \num{22639.1} & \qty{64.61}{\percent} \\
    LP rolling horizon & error $\sigma_{24} = \qty{20}{\percent}$ & \qty{57.82}{\percent} & 0 & \num{22572.9} & \qty{64.42}{\percent} \\
    UB full-year LP & perfect foresight, year & \qty{58.05}{\percent} & 0 & \num{22661.0} & \qty{64.67}{\percent} \\
    \bottomrule
  \end{tabular}
\end{table}

\subsection{Sizing}
\label{sec:eval-sizing}

Each additional server adds fewer heat-feasible hours. With unlimited work, the first server can run
\num{7631} hours a year, the fourth \num{4882}, the sixth \num{3796} and the thirteenth only 5.
Under sourced demand the binding limit is different: at the scaled demand, one to four servers
earn almost the same income (CHF~13,300 to 14,100), because the work is limited by demand in winter
and by heat in summer. One server then has the only positive net after hardware
(Table~\ref{tab:breakeven}): $+$CHF~5,252 a year, against $-$CHF~14,317 for four. Its coverage is
lower (\qty{18.9}{\percent} instead of \qty{22.8}{\percent}). It also misses fewer deadlines
(\qty{5.9}{\percent}) while rejecting more Type~2 work on price (1,250 jobs instead of 878): with
less capacity the load-based price quotes higher, so more summer heat remains for Type~1 work. Four servers would break even at 2.8 times the scaled demand. These results depend on the
demand baseline, a labelled scenario value (Section~\ref{sec:limits}).

\subsection{Local scenario: a pool as a summer heat sink}
\label{sec:eval-local}

A local fleet scenario places five case houses (four servers each, at the Münchenstein electricity
price of CHF~0.3198/kWh) around FHNW Campus Dreispitz and one indoor pool (20 servers) with the
annual heat demand of the renovated Allschwil indoor pool, \qty{627000}{kWh}, spread flat over the
year. It uses sourced, unscaled Basel job volumes, and customers accept every Type~2 price
(elasticity 0). Positions, the flat profile and the pool's 20 servers are scenario values. The same heat-first policy places work where the free share of the hourly
heat budget is largest. In January the pool does 34\% of the work; in July, when the houses only
need hot water, it does 76\%. Over the year 125 of \num{778536} jobs (\qty{0.02}{\percent}) miss
their deadline. The houses cover about 30\% of their heat, the pool 7.9\%. Every site's net after
hardware is still negative at four L4 servers per house (about $-$CHF~12,600 per house).

\subsection{Strategy comparison on a Swiss test network}
\label{sec:eval-network}

The synthetic experiments in this and the next subsections were computed with the research
strategies of Section~\ref{sec:scheduling} at commit \texttt{1863f57}. The test network has 33 sites
in nine Swiss cities:
the reference house, single- and multi-family houses, nine hotels with spa and five indoor pools.
Locations and sizes are \emph{synthetic test data}; heat demand comes from the case formula
(houses) or from the profiles named in the scenario file. All results in this subsection are
\emph{scenarios}, not forecasts of a real network.

At $\varphi=1$, B1 and B2 each cover \qty{83.42}{\percent} of demand with no dumped heat;
B2-reserve covers \qty{83.40}{\percent}. With scarce filler, placement also changes coverage
(Table~\ref{tab:exp-b}). The strategies
differ in missed deadlines (Table~\ref{tab:exp-a}): with \qty{30}{\percent} paid work, B1 misses
\qty{15.69}{\percent} of jobs, while B2 misses none by moving about \num{103000} server-hours to
sites with summer heat demand.

\begin{table}[ht]
  \centering
  \caption{Experiment A: paid-work share at 33 synthetic sites, $\varphi=1$.
    Greedy coverage is \qty{83.42}{\percent} with zero dumped heat
    (B2-reserve: \qty{83.40}{\percent}). The work share is a scenario input.}
  \label{tab:exp-a}
  \small
  \begin{tabular}{rrrrrrr}
    \toprule
    Type 1 & B1 & B2 & B2 moved & B2-reserve & LP & LP moved \\
    share & missed & missed & [server-h] & moved [server-h] & missed & [server-h] \\
    \midrule
    \qty{10}{\percent} & \qty{0}{\percent} & \qty{0}{\percent} & \num{7725} & \num{7725} & \qty{0}{\percent} & \num{0} \\
    \qty{30}{\percent} & \qty{15.69}{\percent} & \qty{0}{\percent} & \num{102754} & \num{64230} & \qty{0}{\percent} & \num{49884} \\
    \qty{50}{\percent} & \qty{25.19}{\percent} & \qty{0}{\percent} & \num{205267} & \num{161968} & \qty{0}{\percent} & \num{149009} \\
    \qty{70}{\percent} & \qty{29.47}{\percent} & \qty{0}{\percent} & \num{309325} & \num{265260} & \qty{0}{\percent} & \num{253300} \\
    \bottomrule
  \end{tabular}
\end{table}

\subsection{Why the network helps}
\label{sec:eval-why}

Figures~\ref{fig:f10} and \ref{fig:f11} show the deadline and routing effects at a glance.
Their underlying scenario values are reported in Tables~\ref{tab:exp-a} and \ref{tab:exp-d}.
% Data from figures/data/a_type1_share.dat, produced by
% python3 scripts/paper_figure_data.py from data/derived/network/experiments/A_type1_share[-lp].csv
\begin{figure}[ht]
  \centering
  \begin{tikzpicture}
    \begin{axis}[name=left, width=0.47\textwidth, height=5.2cm, xlabel={Type~1 share [\%]},
        ylabel={Missed Type~1 deadlines [\%]}, xtick={10,30,50,70}, ymin=0, ymax=35,
        legend style={font=\scriptsize, at={(0.02,0.98)}, anchor=north west}, unbounded coords=jump]
      \addplot+[mark=square*] table[x=share, y=B1missed] {figures/data/a_type1_share.dat};
      \addplot+[mark=triangle*] table[x=share, y=B2missed] {figures/data/a_type1_share.dat};
      \addplot+[mark=*, dashed] table[x=share, y=LPmissed] {figures/data/a_type1_share.dat};
      \legend{B1 local, B2 network, LP}
    \end{axis}
    \begin{axis}[at={(left.east)}, anchor=west, xshift=1.6cm, width=0.47\textwidth, height=5.2cm,
        xlabel={Type~1 share [\%]}, ylabel={Moved Type~1 work [$10^3$ server-h]}, xtick={10,30,50,70},
        ymin=0, legend style={font=\scriptsize, at={(0.02,0.98)}, anchor=north west}, unbounded coords=jump]
      \addplot+[mark=triangle*] table[x=share, y=B2moved] {figures/data/a_type1_share.dat};
      \addplot+[mark=diamond*] table[x=share, y=B2Rmoved] {figures/data/a_type1_share.dat};
      \addplot+[mark=*, dashed] table[x=share, y=LPmoved] {figures/data/a_type1_share.dat};
      \legend{B2, B2-reserve, LP}
    \end{axis}
  \end{tikzpicture}
  \caption{Experiment A at 33 synthetic sites ($\varphi=1$): missed Type~1 jobs (left)
    and moved server-hours (right). Greedy heat coverage is \qty{83.42}{\percent}
    except B2-reserve (\qty{83.40}{\percent}); no heat is dumped.}
  \label{fig:f10}
\end{figure}

% Figure F11 (T-12): data from figures/data/d_network_mix.dat (python3 scripts/paper_figure_data.py)
\begin{figure}[ht]
  \centering
  \begin{tikzpicture}
    \begin{axis}[ybar, bar width=12pt, width=0.6\textwidth, height=5cm, ymin=0, ymax=20,
        ylabel={Missed Type~1 deadlines [\%]}, symbolic x coords={all,residential}, xtick=data,
        xticklabels={{all 33 sites}, {residential only (19)}}, enlarge x limits=0.4,
        legend style={font=\scriptsize, at={(0.02,0.98)}, anchor=north west}, nodes near coords,
        nodes near coords style={font=\scriptsize}, every node near coord/.append style={/pgf/number format/fixed, /pgf/number format/precision=1}]
      \addplot table[x=network, y=B1missed] {figures/data/d_network_mix.dat};
      \addplot table[x=network, y=B2missed] {figures/data/d_network_mix.dat};
      \addplot table[x=network, y=LPmissed] {figures/data/d_network_mix.dat};
      \legend{B1 local, B2 network, LP}
    \end{axis}
  \end{tikzpicture}
  \caption{Experiment D: missed Type~1 jobs at \qty{30}{\percent} paid-work share,
    $\varphi=1$. The synthetic residential-only network has no year-round heat sinks.}
  \label{fig:f11}
\end{figure}


In summer a house can only absorb heat for hot water, about 0.6 servers per day in the case. Paid
work above that cannot run at home without dumping heat. If the network contains year-round heat
sinks (hotels with spas, indoor pools), the work moves there and no deadline is missed. Remove them
and moving still happens, but \qty{13.80}{\percent} of jobs are missed, because every house is short
at the same time (Table~\ref{tab:exp-d}).

\begin{table}[ht]
  \centering
  \caption{Experiment D: network mix at \qty{30}{\percent} Type~1 share and $\varphi=1$.
    All sites are synthetic; B1 and B2 have the same heat coverage.}
  \label{tab:exp-d}
  \small
  \begin{tabular}{lrrrrrrr}
    \toprule
    Network & Sites & B1 & B2 & B2 moved & Coverage & LP & LP moved \\
    & & missed & missed & [server-h] & & missed & [server-h] \\
    \midrule
    all sites & 33 & \qty{15.69}{\percent} & \textbf{\qty{0}{\percent}} & \num{102754} & \qty{83.42}{\percent} & \qty{0}{\percent} & \num{49884} \\
    residential only & 19 & \qty{15.69}{\percent} & \textbf{\qty{13.80}{\percent}} & \num{63333} & \qty{63.18}{\percent} & \qty{13.07}{\percent} & \num{3001} \\
    \bottomrule
  \end{tabular}
\end{table}

Timing at home does not substitute for the network (Table~\ref{tab:reserve}). Keeping tank room for
tomorrow's work (B1-reserve) saves only 0.07 points of missed deadlines (\qty{15.62}{\percent}
versus \qty{15.69}{\percent} for B1), because a daily shortfall returns the next day. Combined with moves (B2-reserve) it reduces the work that has to move by
\qty{37}{\percent}, at a cost of 0.02 points of coverage, because origins keep tank room and need to
send less work away.

\begin{table}[ht]
  \centering
  \caption{Local timing versus moving work.
    Type~1 share \qty{30}{\percent}, $\varphi=1$, 33 synthetic sites.}
  \label{tab:reserve}
  \begin{tabular}{lrrrr}
    \toprule
    Strategy & Coverage & Dumped [kWh] & Missed Type 1 & Moved Type 1 [server-h] \\
    \midrule
    B1 & \qty{83.42}{\percent} & 0 & \qty{15.69}{\percent} & 0 \\
    B1-reserve & \qty{83.42}{\percent} & 0 & \qty{15.62}{\percent} & 0 \\
    B2 & \qty{83.42}{\percent} & 0 & \qty{0}{\percent} & \num{102754} \\
    B2-reserve & \qty{83.40}{\percent} & 0 & \qty{0}{\percent} & \num{64230} \\
    \bottomrule
  \end{tabular}
\end{table}

\paragraph{Transfer delay.} The results above move work without friction. Experiment~E
lets moved work start at the destination only $\delta$ hours after the cohort's
arrival (``data staged at arrival'', a lower bound on friction). The delays are an assumption
sweep; no sourced delay exists. Even $\delta = \qty{12}{h}$ costs no deadline: B2 stays at
\qty{0}{\percent} missed (Table~\ref{tab:exp-e}), because the year-round sinks have enough free
summer capacity to finish the excess in the remaining \qty{12}{h}. B2 even moves less with a delay,
from \num{102754} to \num{76680} and \num{61727} server-hours, i.e.\ from 1.57 to 0.94 moved
server-hours per avoided missed job: without a delay, its persistence test sends work away that would
still have finished at home. With $\delta \ge \qty{24}{h}$ (the deadline) no work can move and missed
deadlines return to B1's level.

\begin{table}[ht]
  \centering
  \caption{Experiment E: assumed transfer delay $\delta$ (33 synthetic sites).
    ``Avoided jobs'' compares each strategy with B1 (\num{65352} missed jobs).
    The LP at $\delta=0$ appears in Table~\ref{tab:exp-a}; $\delta=1$~h was not run with the LP.}
  \label{tab:exp-e}
  \small
  \setlength{\tabcolsep}{4pt}
  \begin{tabular}{rrrrrrr}
    \toprule
    $\delta$ & B2 & B2 moved & B2 moved per & B2-reserve & B2-reserve moved & LP missed \\
    & missed & [server-h] & avoided job & moved [server-h] & per avoided job & / moved \\
    \midrule
    \qty{0}{h} & \qty{0}{\percent} & \num{102754} & 1.57 & \num{64230} & 0.98 & Table~\ref{tab:exp-a} \\
    \qty{1}{h} & \qty{0}{\percent} & \num{102771} & 1.57 & \num{64230} & 0.98 & -- \\
    \qty{4}{h} & \qty{0}{\percent} & \num{76680} & 1.17 & \num{60527} & 0.93 & \qty{0}{\percent} / \num{49793} \\
    \qty{12}{h} & \qty{0}{\percent} & \num{61727} & 0.94 & \num{58775} & 0.90 & \qty{0}{\percent} / \num{49847} \\
    \bottomrule
  \end{tabular}
\end{table}

\subsection{Sensitivity}
\label{sec:eval-sensitivity}

\paragraph{Paid-work volume.} From a Type~1 share of \qty{10}{\percent} to \qty{70}{\percent}, B1's
missed deadlines grow from \qty{0}{\percent} to \qty{29.47}{\percent}; B2 stays at
\qty{0}{\percent}, while the moved server-hours grow roughly linearly (Table~\ref{tab:exp-a}).

\paragraph{Filler availability.} When the filler availability $\varphi$ falls from 1.0 to 0.2,
coverage falls from \qty{83.42}{\percent} to about \qty{21}{\percent} for every strategy
(Table~\ref{tab:exp-b}). A shortage of sellable work, not the algorithm, limits the delivered heat.
B2 helps coverage a little when filler is scarce (0.5 to 0.9 points), because moved Type~1 work
runs where heat is needed and filler would otherwise be missing.

\begin{table}[ht]
  \centering
  \caption{Experiment B: filler availability $\varphi$ at \qty{30}{\percent} Type~1 share
    (33 synthetic sites). The $\varphi=0.35$ scenario is the filler-capacity analogy of the demand
    research in the repository (\texttt{docs/research/demand.md}).}
  \label{tab:exp-b}
  \begin{tabular}{rrrrr}
    \toprule
    $\varphi$ & B1 coverage & B2 coverage & B1 missed & B2 missed \\
    \midrule
    1.0 & \qty{83.42}{\percent} & \qty{83.42}{\percent} & \qty{15.69}{\percent} & \qty{0}{\percent} \\
    0.6 & \qty{54.21}{\percent} & \qty{54.71}{\percent} & \qty{15.54}{\percent} & \qty{0}{\percent} \\
    0.35 & \qty{34.58}{\percent} & \qty{35.38}{\percent} & \qty{14.93}{\percent} & \qty{0}{\percent} \\
    0.2 & \qty{21.41}{\percent} & \qty{22.34}{\percent} & \qty{14.82}{\percent} & \qty{0}{\percent} \\
    \bottomrule
  \end{tabular}
\end{table}

\paragraph{Forecast error in the network.} In the synthetic error sweep, LP coverage falls
from \qty{83.42}{\percent} to \qty{83.23}{\percent} at
$\sigma_{24}=\qty{20}{\percent}$, with no missed deadlines or dumped heat
(Table~\ref{tab:exp-c}). This tests the stated error process rather than a measured forecast.

\begin{table}[ht]
  \centering
  \caption{Experiment C: LP forecast-error sensitivity (33 synthetic sites,
    \qty{30}{\percent} Type~1 share, $\varphi=1$). AR(1) parameters $\rho=0.8$ and
    $\sigma_{24}$ are uncalibrated; the perfect-foresight row is from experiment~A.}
  \label{tab:exp-c}
  \begin{tabular}{lrrrr}
    \toprule
    Forecast & Coverage & Dumped [kWh] & Missed Type 1 & Moved [server-h] \\
    \midrule
    perfect foresight & \qty{83.42}{\percent} & 0 & \qty{0}{\percent} & \num{49884} \\
    $\sigma_{24} = 0.05$ & \qty{83.39}{\percent} & 0 & \qty{0}{\percent} & \num{49882} \\
    $\sigma_{24} = 0.10$ & \qty{83.35}{\percent} & 0 & \qty{0}{\percent} & \num{49929} \\
    $\sigma_{24} = 0.20$ & \qty{83.23}{\percent} & 0 & \qty{0}{\percent} & \num{49967} \\
    \bottomrule
  \end{tabular}
\end{table}

\paragraph{Economic sensitivities.} Compute price, hardware, electricity price and heat fee: see
Section~\ref{sec:economics}.

\subsection{Summary of findings}

\begin{enumerate}
  \item[(i)] With unlimited work four servers cover 58\% of the reference house's heat; with sourced
    demand they cover 22.8\%, dump nothing and miss 8.9\% of deadlines, mostly in summer.
  \item[(ii)] At the sourced demand, hardware decides the result: one server covers its cost,
    four need 2.8 times the demand or 2.6 times the Type~1 price.
  \item[(iii)] Assuming 86\% boiler efficiency, the case formula fits a real house's space heating; the real house's larger summer
    heat use removes most of the summer deadline misses.
  \item[(iv)] On one house the greedy rule reaches the full-year upper bound. Value from routing
    arises in a network with year-round heat sinks: the executable policy moves 76\% of July work
    to the Dreispitz pool, and in the synthetic experiments greedy routing (B2) already removes the
    missed deadlines; the LP mainly reduces the work moved.
  \item[(v)] The limits of these statements are listed in Section~\ref{sec:limits}.
\end{enumerate}

\clearpage % Keep evaluation tables before the next section.
\section{Limits and threats to validity}
\label{sec:limits}

\paragraph{Physical model.} The case-defined reference house assumes a lossless tank, 80\% heat
recovery, constant daily hot-water demand, flat electricity price and abstract \qty{1}{kW}
servers. The smaller four-L4 variant uses host-power estimates and GPU TDP rather than measured
wall power or a quoted liquid-cooled build. Its calculated heat fits the constant hot-water
base load; actual draw timing, delivery temperature, thermal losses and water hygiene require
engineering validation. Neither normalised 1~kW L4 profile nor the candidate four-card host is
an installed bill of materials.

\paragraph{Work and markets.} The reference-house maximum assumes unlimited paid work. The
mixed-queue run uses a medical-dictation arrival fixture for technical validation, not the
proposed public-document LLM market; sensitive health records are not placed in homes. Its
0.35 filler capacity cap is an analogy, not observed sales. The laptop Apertus measurement is
CPU-package energy, not L4 throughput or whole-node power. The research simulation was measured
on a CPU laptop; a billable GPU-compatible filler has not been demonstrated. Competitor list
prices and retail GPU rates are buyer prices, not DigitalHeat contracts or net receipts. GPU wear
and purchase prices are sensitivity inputs. Cooling integration, idle energy, connectivity,
maintenance, insurance, financing, tax bases and sales costs remain incompletely priced. The
consumer-GPU cost case is excluded from the proposed deployment pending licence clarification.

\paragraph{Network inference.} The 33-site locations and sizes and Type~1 arrival shares are
synthetic. Unless a forecast-error sweep is explicitly stated, planning uses perfect heat
foresight while knowing only arrived jobs. Uniform transfer delay assumes data staged at arrival;
transfer cost and distance-dependent delay are absent. The experiments therefore compare
algorithms under controlled inputs rather than predict an operating Swiss network.

\section{Conclusion}
\label{sec:conclusion}

The case reference house sets a physical limit: with unlimited work, four abstract servers could
deliver \qty{18129}{kWh} of useful heat, 58\% of the demand, and a simple local controller already
reaches the full-year upper bound. With finite, sourced job demand the picture changes. The same
four servers deliver \qty{7133}{kWh} (22.8\%), dump nothing, and miss 8.9\% of deadlines, nearly
all for work that arrives in summer when the house only needs hot water. Demand, not heat, limits
the house in winter; heat limits it in summer.

Economic viability is an actor-level question. Income exceeds the electricity cost, but at the sourced demand four GPU servers
lose about CHF~14,300 a year after hardware. One server earns almost the same income and is the
only configuration that covers its hardware: the combined system keeps about CHF~5,250 a year, the
operator alone, who gives the heat away, about CHF~4,470 (payback 3.4 years). Four servers need 2.8 times the demand or 2.6 times
the Type~1 price; raising the research price alone does not suffice. The sizing question of the
case therefore has a practical answer under these assumptions: size for the work that can be sold,
not for the winter peak.

Assuming 86\% boiler efficiency, the case formula describes a real house's space heating within
1.3\% of the case's design power,
while that house's summer heat use is three times the case's hot water and removes most summer
deadline misses. Year-round heat sinks serve the same role in a network: a nearby pool takes 76\%
of the July work, and in synthetic experiments with earlier research strategies, routing reduces missed deadlines where houses alone
cannot. The demonstrator shows the scheduler freezing and resuming a real container on the heat
signal, with measured CPU time and energy.

The next tests are a measured, thermally integrated node, a customer commitment stating price and
minimum paid capacity, and a measured hot-water profile; the thermal test must record delivery
temperature, circulation-pump power and storage losses so that measured useful heat can replace
the idealised recovery assumption. The paper's contribution is a reproducible boundary between
demonstrated heat control, conditional home economics and network scenarios.


\appendix
\section{Reproducibility}
\label{app:repro}
The case weather and price inputs have \texttt{SOURCE} files in \texttt{data/raw/}; every
configuration value carries its source, and \texttt{SOURCES.md} lists every external source with
the place where it is used. Fleet runs are deterministic (fixed seeds) and write their resolved
scenario, seeds and input hashes to \texttt{run.json}. Fleet outputs are not committed; the
commands below regenerate them in about a second per house. The specification, every input of the
reference run with its source, and the known simplifications are in
\texttt{docs/model/fleet-simulation.md}. The synthetic network experiments use research strategies
that are not selectable on the current main branch; they are reproducible at commit
\texttt{1863f57}, which also holds their output files in \texttt{data/derived/network/experiments/}.

\begin{table}[ht]
  \centering
  \caption{One command per result. Paths are relative to the repository root; ``Fleet'' stands for
    \texttt{dotnet run -c Release -{}-project src/DigitalHeat.FleetRun -{}-}.}
  \label{tab:repro}
  \begin{tabular}{>{\raggedright\arraybackslash}p{0.28\linewidth}>{\raggedright\arraybackslash}p{0.66\linewidth}}
    \toprule
    Result & Command and output \\
    \midrule
    Table~\ref{tab:house-fleet}, reference house &
      Fleet \texttt{run config/fleet/scenarios/reference.json}, then Fleet \texttt{post data/derived/fleet/reference} (hardware) \\
    Table~\ref{tab:house-fleet}, unlimited work &
      \texttt{dotnet run -{}-project src/DigitalHeat.YearRun -{}- .}; fleet cross-check:
      \texttt{dotnet test src/DigitalHeat.Fleet.Tests -{}-filter ReferenceHouseWithBacklogMatchesCore} \\
    Table~\ref{tab:house-fleet}, real house &
      \texttt{python3 -I scripts/calibrate\_real\_house.py} (needs the organisers' data, see
      \texttt{data/raw/housedata/SOURCE.md}), then Fleet \texttt{run config/fleet/scenarios/real-house.json} \\
    Table~\ref{tab:breakeven} &
      copies of \texttt{reference.json} with one change each (\texttt{servers.n},
      \texttt{demand.volume\_scale}, or the job types' \texttt{price.chf\_per\_work\_unit}), each run and
      post-processed as above; values in \texttt{docs/model/fleet-simulation.md} \S16 \\
    Local scenario (Section~\ref{sec:eval-local}) &
      Fleet \texttt{run config/fleet/scenarios/dreispitz-5-pool.json} \\
    Table~\ref{tab:house-strategies} (one house, research strategies) &
      at commit \texttt{1863f57}: \texttt{dotnet run -{}-project src/DigitalHeat.NetworkRun -{}- .} \\
    Tables~\ref{tab:exp-a}, \ref{tab:exp-d}, \ref{tab:reserve}, \ref{tab:exp-e}, \ref{tab:exp-b}, \ref{tab:exp-c} (network, research strategies) &
      at commit \texttt{1863f57}: \texttt{python3 scripts/network\_experiments.py} (greedy) and
      \texttt{-{}-lp -{}-jobs 4} (LP) $\to$ \texttt{data/derived/network/experiments/*.csv};
      each CSV row records its command line \\
    Network figures &
      \texttt{python3 scripts/paper\_figure\_data.py} $\to$ \texttt{docs/paper/latex/figures/data/*.dat} (runs on current main from the committed CSVs) \\
    Unlimited-work sizing (Section~\ref{sec:eval-sizing}) &
      \texttt{dotnet run -{}-project src/DigitalHeat.YearRun -{}- . -{}-sizing data/derived/server-sizing/reference-house.json} \\
    Operator and small-node sensitivities (Section~\ref{sec:economics}) &
      Worked arithmetic and input ledgers in \texttt{docs/research/free-heat-viability.md} and
      \texttt{docs/research/profitable-home-path.md}; no full automated rerun yet \\
    Market price screen (Section~\ref{sec:economics}) &
      \texttt{python3 scripts/price\_comparison.py}; inputs and limits in
      \texttt{docs/research/price-comparison.md} \\
    Demonstrator (Section~\ref{sec:demonstrator}) &
      \texttt{dotnet run -c Release -{}-project src/DigitalHeat.VoiceDemo}; runbook, freeze probe and
      measured run in \texttt{docs/demo/voice-request.md} \\
    Measured laptop jobs (Section~\ref{sec:demonstrator}) &
      Logs and commands in \texttt{docs/jobs/llm-summary.md} and
      \texttt{docs/jobs/research-sim.md} \\
    Tests &
      \texttt{dotnet test src/DigitalHeat.sln} \\
    \bottomrule
  \end{tabular}
\end{table}

\clearpage % Keep the reproducibility table ahead of the references.

\bibliographystyle{unsrt}
\bibliography{references}

\end{document}
